% ==============================================================================
% Benchmark
% ==============================================================================

\section{Benchmark}
\label{sec:benchmark}

We construct a \emph{contrastive} cognitive bias benchmark designed to elicit
and control for the S1-like/S2-like processing distinction.  The design follows
three principles: (1)~every conflict item has a structurally matched no-conflict
control \citep{deneys2012bias}; (2)~all items use novel surface features to
resist memorization-based shortcuts \citep{hagendorff2023human}; and (3)~a $2
\times 2$ difficulty $\times$ congruence design disentangles processing mode
from task difficulty.

\subsection{Task Categories}

The benchmark spans seven categories, each targeting a well-characterized
heuristic-deliberative dissociation from the cognitive science literature:

\begin{enumerate}[leftmargin=*,itemsep=1pt]
  \item \textbf{CRT variants} (48 conflict / 48 no-conflict).  Structural
  isomorphs of the Cognitive Reflection Test with modified surface features
  (changed objects, domains, and numerical values).  We draw on the 150
  decontaminated items from \citet{hagendorff2023human}'s OSF repository and
  generate additional structural variants via controlled template perturbation.

  \item \textbf{Base rate neglect} (22 / 22).  Bayesian reasoning problems
  where a diagnostic description cues a low-base-rate category.  Conflict
  items set the base rate against the description; no-conflict items align
  them.  All posterior probabilities are analytically verified.

  \item \textbf{Belief-bias syllogisms} (20 / 20).  Syllogisms crossing
  logical validity ($\times 2$) with conclusion believability ($\times 2$),
  following the Evans taxonomy.  Conflict items have valid--unbelievable or
  invalid--believable conclusions; controls have congruent
  validity--believability.

  \item \textbf{Anchoring} (14 / 14).  Numerical estimation tasks preceded by
  a high or low anchor.  Conflict items use extreme anchors designed to
  distort estimates; no-conflict items use anchors near the correct value.

  \item \textbf{Framing effects} (14 / 14).  Gain- versus loss-framed
  decision problems.  Conflict items present logically equivalent options in
  different frames; controls present options where framing and value align.

  \item \textbf{Conjunction fallacy} (12 / 12).  Probability judgment tasks
  where a detailed description makes the conjunction of two events seem more
  likely than a single constituent.  Novel scenarios replace the classic
  ``Linda problem.''

  \item \textbf{Multi-step arithmetic} (12 / 12).  Graded-difficulty
  arithmetic (1--5 operations) with and without carrying traps.  Conflict
  items embed a salient but incorrect shortcut; no-conflict items do not.
\end{enumerate}

\subsection{Benchmark Statistics}

The final benchmark comprises \textbf{284 items} organized as \textbf{142
matched pairs} (conflict + no-conflict).  Each item includes 2--5
paraphrases for a total of \textbf{558 unique prompts}, allowing us to
assess robustness to surface variation.  Items are presented in two formats:
constrained free-response (``Answer with just a number:'') for activation
analysis, and multiple-choice for behavioral validation.

\paragraph{Scoring.}  We use a three-tier scoring system: (1)~\emph{binary}
(correct/incorrect); (2)~\emph{categorical} (correct / S1-lure / other-wrong
/ refusal)---the primary target for probing; and (3)~\emph{continuous}
(distance from the S2-correct vs.\ S1-predicted answer).

\paragraph{Contamination controls.}  We include classic problems (e.g., the
original bat-and-ball) as contamination baselines only---they are excluded from
all primary analyses.  We test memorization by prompting models to complete the
original CRT verbatim and comparing performance on classic versus novel variants.

\subsection{Example Items}

Table~\ref{tab:examples} illustrates the conflict/no-conflict matched-pair
structure for three representative categories.

\begin{table}[t]
\centering
\caption{Example conflict/no-conflict item pairs from three benchmark
  categories.  Each pair shares the same correct answer; they differ only in
  whether a heuristic lure is present.}
\label{tab:examples}
\vspace{4pt}
\small
\begin{tabular}{p{1.2cm}p{3.0cm}p{3.0cm}cc}
\toprule
\textbf{Category} & \textbf{Conflict item} & \textbf{No-conflict control}
  & \textbf{Lure} & \textbf{Correct} \\
\midrule
CRT &
  A notebook and a pen cost \$2.20 together.  The notebook costs \$2.00
  more than the pen.  How much does the pen cost? &
  A notebook and a pen cost \$2.10 together.  The notebook costs \$2.00.
  How much does the pen cost? &
  \$0.20 & \$0.10 \\[6pt]
Base rate &
  In a group of 1000 people, 5 are engineers and 995 are lawyers.  Tom
  likes puzzles and circuits.  Is Tom more likely an engineer or a
  lawyer? &
  In a group of 1000 people, 500 are engineers and 500 are lawyers.  Tom
  likes puzzles and circuits.  Is Tom more likely an engineer or a
  lawyer? &
  Engineer & Lawyer \\[6pt]
Syllogism &
  All flowers need water.  Roses need water.  Therefore roses are
  flowers.  Is this valid? &
  All mammals are warm-blooded.  Dogs are mammals.  Therefore dogs are
  warm-blooded.  Is this valid? &
  Valid & Invalid \\
\bottomrule
\end{tabular}
\end{table}

\subsection{Scoring Detail}

We employ three scoring tiers, applied automatically via regex patterns
stored alongside each benchmark item:
\begin{enumerate}[leftmargin=*,itemsep=1pt]
  \item \textbf{Binary}: 1 if the model's response matches the
  \texttt{correct\_answer} regex, 0 otherwise.
  \item \textbf{Categorical}: one of \{correct, S1-lure, other-wrong,
  refusal\}.  The S1-lure label is assigned when the response matches the
  \texttt{lure\_answer} regex (conflict items only).  Refusals are detected
  via keyword matching.  This four-way classification is the primary probing
  target.
  \item \textbf{Continuous} (numerical categories only):
  $|r - c| / |l - c|$, where $r$ is the extracted numerical response, $c$
  is the correct answer, and $l$ is the lure answer.  Values near 0
  indicate correctness; values near 1 indicate the lure response; values
  $> 1$ indicate errors worse than the lure.
\end{enumerate}

\subsection{Behavioral Validation}

\textbf{[Placeholder: behavioral results on all four models.]}  We report the
proportion of S1-lure responses on conflict versus no-conflict items, broken
down by category and model.  The benchmark is validated if (i)~models produce
$>30$\% lure responses on conflict items and (ii)~standard models show
significantly higher lure rates than reasoning-trained models.
% Expected: Table 1 with behavioral accuracy by model x category x condition
